\clearpage\SourceLeaf{451}%
\fancyhead[L]{\scriptsize BEMERKUNGEN.}
\pdfbookmark[1]{BEMERKUNGEN.}{bemerkungen-coordinaten}
\begin{center}\textbf{BEMERKUNGEN.}\end{center}

\indent Der Einheit der Coordinaten so wie den verschiedenen Reductionen der Messung sollten vermuthlich die von \textsc{Walbeck} gefundenen Erddimensionen zu Grunde gelegt werden.

\textsc{Walbeck} et \textsc{Brummer}. De forma et magnitudine telluris. Aboae 1819 pag. 16: ‘Gradus medius seu \(\frac{1}{90}\) pars Quadrantis Meridiani \(=57009^{t}{,}76\). Ellipticitas \(=\frac{1}{302{,}78}\)’ [Handschriftliche Bemerkung von \textsc{Gauss}: mittlere Meridiangrad] \(=\) ‘\(57009^{t}{,}7584\). Der Meter also \(=443^{l}{,}307885\), Verhältniss \(=37299:37300\) Logarithm \(=0{,}00001164\).’

\textsc{Gauss}. Bestimmung des Breitenunterschiedes zwischen den Sternwarten von Göttingen und Altona. Göttingen 1828. Art. 20. — ‘Wenn man meine Dreiecke als auf der Oberfläche eines elliptischen Sphaeroids liegend, dessen Dimensionen die von \textsc{Walbeck} aus der Gesammtheit der bisherigen Gradmessungen abgeleiteten sind, und welches nach unsrer besten gegenwärtigen Kenntniss sich am vollkommensten an die wirkliche Gestalt \textit{im Ganzen} anschliesst (Abplattung \(\frac{1}{302{,}76}\), der dreihundertsechzigste Theil des Erdmeridians \(=57009{,}746\) Toisen) berechnet, und dabei von der Polhöhe von Göttingen \(=51^{0}\ 31'\ 37''85\) ausgeht’..

Hienach scheint \textsc{Gauss} mehrfach mit der Abplattung \(\frac{1}{302{,}76}\) statt mit der \textsc{Walbeck}schen \(\frac{1}{302{,}78}\) gerechnet zu haben und in der That liegt auch mehren der noch im handschriftlichen Nachlass vorhandenen Hülfstafeln die erstere Zahl zu Grunde.

\textsc{Gauss} an \textsc{Schumacher}. Göttingen 1830 April 18 ‘Zweite Hülfstafel, Anmerkung: ‘Bei früher von mir mitgetheilten Coordinaten ist die Einheit \(\frac{1}{10000000}\) des Erdquadranten nach \textsc{Walbeck}’s Dimensionen; um jene also in solche zu verwandeln, bei denen die Einheit \(\frac{1}{10000000}\) des Erdquadranten nach \textsc{Schmidt}’s neuesten zum Grunde liegt, müssen jene erst mit \(\frac{57009758}{57008551}\) oder mit \(1+\frac{1}{47245}\) multiplicirt werden.’

\textsc{Gauss}. Bestimmung des Breitenunterschiedes zwischen Göttingen und Altona. Art. 19. ‘Nach der trigonometrischen Verbindung der Sternwarten von Göttingen und Altona liegt letztere 115163,725 Toisen nördlich, 7,211 Toisen westlich von jener. Diese Zahlen beziehen sich auf die Plätze der Meridiankreise; sie gründen sich auf den Werth der Dreiecksseite Hamburg-Hohenhorn 13841,815 Toisen, und diese auf die von Hrn. Prof. \textsc{Schumacher} in Holstein im Jahre 1820 gemessenen Basis. Da jedoch die Vergleichung der dabei gebrauchten Messstangen mit der Normaltoise noch nicht \textit{definitiv} vollendet ist, so wird obige Entfernung in Zukunft noch in demselben Verhältniss abzuändern sein, wie die Basis selbst, welche Veränderung aber jedenfalls nur sehr gering sein kann.’

Herr Geheimer Etatsrath \textsc{Andrae} in Copenhagen bemerkt über die Revision der Basis in einem Schreiben vom 5 März 1865 abgedruckt im Generalbericht über die mitteleuropäische Gradmessung für das Jahr 1864 Seite 6. 7. ‘Die von \textsc{Schumacher} angegebene Länge der \textit{Braacker Basis}: 3014,5799 Toisen, welche bei den früheren Berechnungen sowohl der Dänischen als auch der Hannöverschen unter der Leitung von \textsc{Gauss} ausgeführten Triangulationen angewendet wurde, konnte nur als ein vorläufiges Resultat der Basismessung angesehen werden, da die Reduction auf den Meeresspiegel und mehrere andere Correctionen noch nicht berücksichtigt waren. Da diese Reductionen an Grösse beträchtlich die Unsicherheiten der Messungen selbst, die mit grosser Sorgfalt ausgeführt sind, übersteigen, war eine neue Bestimmung nothwendig und Herr Professor Dr. \textsc{Peters} in Altona hat auch die Güte gehabt, eine ausführliche, mit der grössten Genauigkeit durchgeführte Berechnung sämmtlicher Correctionen vorzunehmen, durch welche die Länge der Basis sich nun stellt wie folgt:
